%%%%%%%%%%%%%%%%%%%%%%%%%%%%%%%%%%%%%%%%%%%%%%%%%%%%%%%%%%%%%%%%%%%%%%%%%%%%%%%
%  Kapitel 1 -- Einleitung
%
%  Dieses Kapitel ist als Geruest gedacht: Die Zwischenueberschriften decken
%  ab, was eine Einleitung ueblicherweise leistet. Text ersetzen, nicht
%  benoetigte Abschnitte loeschen.
%
%  Bilder liegen in Pictures/ und werden OHNE diesen Ordnernamen und OHNE
%  Dateiendung angegeben -- \graphicspath erledigt den Rest.
%%%%%%%%%%%%%%%%%%%%%%%%%%%%%%%%%%%%%%%%%%%%%%%%%%%%%%%%%%%%%%%%%%%%%%%%%%%%%%%

\chapterimage{chapter_image_1.jpg} % Bild im Kapitelkopf (Seitenverhaeltnis 2:1)

\chapter{Einleitung}
\label{ch:einleitung}
\index{Einleitung}


%==============================================================================
\section{Motivation}
\label{sec:motivation}
%==============================================================================

Hier steht, warum das Thema überhaupt behandelt wird: das Problem, die
Ausgangslage, der Anlass. Anführungszeichen setzt man mit \enquote{enquote},
dann stimmen sie automatisch in jeder Sprache -- verschachtelt funktioniert
das ebenfalls, etwa \enquote{ein Zitat mit \enquote{Zitat darin}}.

Zahlen mit Einheiten gehören in \verb|\SI|, damit Abstand und Dezimalzeichen
einheitlich bleiben: Die Messung ergab \SI{12345.67}{\metre\per\second} bei
einer Spannung von \SI{3.3}{\volt}. Im Quelltext schreibt man dabei immer
einen Punkt -- die Ausgabe als Komma erledigt \texttt{siunitx}.

\lipsum[1] % Blindtext -- vor der Abgabe entfernen


%==============================================================================
\section{Zielsetzung}
\label{sec:zielsetzung}
%==============================================================================

Was soll am Ende erreicht sein? Idealerweise so konkret, dass sich im Fazit
überprüfen lässt, ob es erreicht wurde.

Auf andere Stellen verweist man mit \verb|\cref|. Das Wort davor liefert der
Befehl bereits mit -- man schreibt also nur \verb|\cref{ch:einleitung}| und
erhält \cref{ch:einleitung}. Ein zusätzliches \enquote{Kapitel} davor wäre
doppelt. Für Abschnitte, Abbildungen und Tabellen gilt dasselbe:
\cref{sec:motivation} und \cref{fig:beispielbild}.

\ofig[beispielbild]{Ein Platzhalterbild. Die Bildunterschrift steht bei
Abbildungen unter dem Bild und endet ohne Punkt}{placeholder}{0.6\linewidth}

Abkürzungen werden bei der ersten Verwendung automatisch ausgeschrieben:
\acrfull{gcd}. Ab dann genügt die Kurzform \acrshort{gcd}. Das
Abkürzungsverzeichnis am Ende des Dokuments füllt sich von selbst -- neue
Einträge kommen in \texttt{acronyms.tex}.


%==============================================================================
\section{Aufbau der Arbeit}
\label{sec:aufbau}
%==============================================================================

Ein kurzer Wegweiser durch die folgenden Kapitel. Quellen werden mit
\verb|\cite| zitiert \cite{book_key}; mit Seitenangabe sieht das so
aus \cite[162]{book_key}. Die Einträge dafür stehen in der Datei
\texttt{bibliography.bib}.

\lipsum[2] % Blindtext -- vor der Abgabe entfernen
